% !TeX program = xelatex
% !TeX encoding = UTF-8
% !TeX spellcheck = pt_BR
% ----------------------------------------------------------------------
% Guia Prático para a Capacitação de MySQL
% Compilar com XeLaTeX (a fonte Lexend é carregada localmente):
%     latexmk -xelatex guia.tex
% ----------------------------------------------------------------------

\documentclass[12pt, a4paper, openright, oneside, brazil]{abntex2}

% ----------------------------------------------------------------------
% Preâmbulo — Capacitação de MySQL
% ----------------------------------------------------------------------
% ATENÇÃO: compilar com XeLaTeX. A fonte Lexend é carregada de arquivos
% locais em assets/fonts/, o que exige fontspec (XeLaTeX ou LuaLaTeX).
% ----------------------------------------------------------------------
\usepackage[brazil]{babel}

% --- Fonte da marca (Lexend, local) -----------------------------------
\usepackage{fontspec}
\newcommand{\fontpath}{./assets/fonts/}
\setmainfont[
    Path          = \fontpath,
    UprightFont   = *-Regular,
    BoldFont      = *-Bold,
    AutoFakeSlant = 0.2,          % Lexend não possui itálico real
]{Lexend}

% --- Estrutura e tipografia -------------------------------------------
\usepackage{microtype}
\usepackage{indentfirst}
\usepackage{graphicx}
\usepackage{xcolor}
\usepackage{float}
\usepackage{fancyhdr}

% --- Tabelas -----------------------------------------------------------
\usepackage{longtable}
\usepackage{tabularx}
\usepackage{booktabs}
\usepackage{array}
\usepackage{multirow}

% --- Código ------------------------------------------------------------
\usepackage{listings}

% --- Símbolos usados nos avisos ---------------------------------------
\usepackage{amssymb}
\usepackage{pifont}

% --- Caixas de destaque ------------------------------------------------
\usepackage[most]{tcolorbox}

% --- Links e citações --------------------------------------------------
\usepackage{url}
\usepackage{hyperref}
\usepackage{abntex2cite}

% --- Justificação -----------------------------------------------------
% Nomes de comandos em \texttt{} são longos e não hifenizam, o que empurra
% linhas para fora da margem. O emergencystretch autoriza o TeX a esticar
% mais os espaços nesses parágrafos em vez de estourar a caixa.
\setlength{\emergencystretch}{3em}

% ----------------------------------------------------------------------
% Configurações visuais — Capacitação de MySQL
% ----------------------------------------------------------------------

% --- Paleta da marca ---------------------------------------------------
\definecolor{branco}{HTML}{FFFFFF}
\definecolor{roxoescuro}{HTML}{1E1647}
\definecolor{roxoclaro}{HTML}{614AD3}
\definecolor{amarelo}{HTML}{FFFF00}
\definecolor{cinzafundo}{HTML}{F4F4F7}
\definecolor{cinzalinha}{HTML}{9A9AB0}
\definecolor{verdeok}{HTML}{1B7F4B}
\definecolor{vermelhoerro}{HTML}{B3261E}

% --- Atalhos de marcação -----------------------------------------------
\newcommand{\comando}[1]{\texttt{#1}}
\newcommand{\sqlkw}[1]{\texttt{\textbf{#1}}}
\newcommand{\ok}{\textcolor{verdeok}{\ding{51}}}
\newcommand{\nok}{\textcolor{vermelhoerro}{\ding{55}}}

% ----------------------------------------------------------------------
% Ambientes de código (listings)
% ----------------------------------------------------------------------
% O 'literate' abaixo é necessário porque o listings não processa UTF-8
% sozinho: sem ele, cada acento dentro de um bloco de código vira lixo.
\lstset{
    inputencoding=utf8,
    extendedchars=true,
    literate=%
      {á}{{\'a}}1 {à}{{\`a}}1 {ã}{{\~a}}1 {â}{{\^a}}1
      {é}{{\'e}}1 {ê}{{\^e}}1 {í}{{\'i}}1
      {ó}{{\'o}}1 {õ}{{\~o}}1 {ô}{{\^o}}1
      {ú}{{\'u}}1 {ü}{{\"u}}1 {ç}{{\c c}}1
      {Á}{{\'A}}1 {À}{{\`A}}1 {Ã}{{\~A}}1 {Â}{{\^A}}1
      {É}{{\'E}}1 {Ê}{{\^E}}1 {Í}{{\'I}}1
      {Ó}{{\'O}}1 {Õ}{{\~O}}1 {Ô}{{\^O}}1
      {Ú}{{\'U}}1 {Ç}{{\c C}}1
      {º}{{\textordmasculine}}1 {ª}{{\textordfeminine}}1
}

\lstdefinestyle{base}{
    basicstyle       = \ttfamily\footnotesize,
    backgroundcolor  = \color{cinzafundo},
    frame            = single,
    rulecolor        = \color{cinzalinha},
    breaklines       = true,
    breakatwhitespace= false,
    postbreak        = \mbox{\textcolor{roxoclaro}{$\hookrightarrow$}\space},
    showstringspaces = false,
    upquote          = true,
    tabsize          = 2,
    captionpos       = b,
    xleftmargin      = 12pt,
    framexleftmargin = 12pt,
}

\lstdefinestyle{sql}{
    style        = base,
    language     = SQL,
    morekeywords = {AUTO_INCREMENT,ENGINE,InnoDB,UNSIGNED,IF,EXISTS,DATABASE,
                    REFERENCES,CONSTRAINT,CASCADE,RESTRICT,BOOLEAN,TINYINT,
                    DECIMAL,VARCHAR,TIMESTAMP,CHARACTER,COLLATE,LIMIT,SHOW,
                    DESCRIBE,TRUNCATE,USE,REPLACE},
    keywordstyle = \color{roxoescuro}\bfseries,
    commentstyle = \color{cinzalinha}\itshape,
    stringstyle  = \color{roxoclaro},
}

\lstdefinestyle{bash}{
    style        = base,
    language     = bash,
    keywordstyle = \color{roxoescuro}\bfseries,
    commentstyle = \color{cinzalinha}\itshape,
    stringstyle  = \color{roxoclaro},
}

\lstdefinestyle{yaml}{
    style         = base,
    basicstyle    = \ttfamily\footnotesize\color{roxoescuro},
    commentstyle  = \color{cinzalinha}\itshape,
    morecomment   = [l]{\#},
    moredelim     = [l][\color{black}]{:\ },
}

\lstnewenvironment{sqlcode}[1][]{\lstset{style=sql,#1}}{}
\lstnewenvironment{bashcode}[1][]{\lstset{style=bash,#1}}{}
\lstnewenvironment{yamlcode}[1][]{\lstset{style=yaml,#1}}{}

% ----------------------------------------------------------------------
% Caixas de destaque
% ----------------------------------------------------------------------
\newtcolorbox{aviso}[1][Atenção]{
    enhanced, breakable,
    colback   = vermelhoerro!5,
    colframe  = vermelhoerro,
    coltitle  = branco,
    fonttitle = \bfseries,
    title     = {#1},
    left=3mm, right=3mm, top=2mm, bottom=2mm,
}

\newtcolorbox{dica}[1][Dica]{
    enhanced, breakable,
    colback   = roxoclaro!5,
    colframe  = roxoclaro,
    coltitle  = branco,
    fonttitle = \bfseries,
    title     = {#1},
    left=3mm, right=3mm, top=2mm, bottom=2mm,
}

\newtcolorbox{conceito}[1][Conceito]{
    enhanced, breakable,
    colback   = roxoescuro!5,
    colframe  = roxoescuro,
    coltitle  = branco,
    fonttitle = \bfseries,
    title     = {#1},
    left=3mm, right=3mm, top=2mm, bottom=2mm,
}

% ----------------------------------------------------------------------
% Links e metadados do PDF
% ----------------------------------------------------------------------
\hypersetup{
    colorlinks  = true,
    linkcolor   = roxoescuro,
    citecolor   = roxoescuro,
    filecolor   = roxoclaro,
    urlcolor    = roxoclaro,
    pdftitle    = {Guia Prático para a Capacitação de MySQL},
    pdfauthor   = {Ronivaldo Domingues de Andrade},
    pdfsubject  = {Fundamentos de MySQL e Projetos de Bancos de Dados},
    pdfkeywords = {MySQL, SQL, Banco de Dados, SGBD, Modelo Relacional,
                   CRUD, JOIN, phpMyAdmin, Docker},
}


% --- Informações da capa e da folha de rosto --------------------------
\titulo{Fundamentos do MySQL e Projetos de Bancos de Dados}
\autor{Ronivaldo Domingues de Andrade}
\local{Rio de Janeiro - RJ}
\data{2026}
\tipotrabalho{Guia Prático para Capacitação}
\preambulo{Guia prático da capacitação de MySQL e SQL, com duração de duas
horas, elaborado para os membros da liga acadêmica.}

\begin{document}
\selectlanguage{brazil}

% --- Elementos pré-textuais -------------------------------------------
\imprimircapa
\imprimirfolhaderosto

\newpage
\listoftables

\begin{siglas}
  \item[SGBD] Sistema de Gerenciamento de Banco de Dados
  \item[DBMS] \emph{Database Management System} (o mesmo que SGBD, em inglês)
  \item[RDBMS] \emph{Relational Database Management System}, SGBD relacional
  \item[SQL] \emph{Structured Query Language}, linguagem de consulta estruturada
  \item[DDL] \emph{Data Definition Language}, subconjunto do SQL que define estruturas
  \item[DML] \emph{Data Manipulation Language}, subconjunto do SQL que manipula dados
  \item[DQL] \emph{Data Query Language}, subconjunto do SQL que consulta dados
  \item[DCL] \emph{Data Control Language}, subconjunto do SQL que controla permissões
  \item[TCL] \emph{Transaction Control Language}, subconjunto do SQL que controla transações
  \item[CRUD] \emph{Create, Read, Update, Delete}, as quatro operações básicas sobre dados
  \item[PK] \emph{Primary Key}, chave primária
  \item[FK] \emph{Foreign Key}, chave estrangeira
  \item[InnoDB] Motor de armazenamento padrão do MySQL, com suporte a chaves estrangeiras
  \item[LTS] \emph{Long Term Support}, versão com suporte de longo prazo
  \item[ORM] \emph{Object-Relational Mapping}, mapeamento objeto-relacional
\end{siglas}

\tableofcontents

% --- Capítulos ---------------------------------------------------------
\chapter{Introdução}

Todo sistema que você vai construir na liga --- um cadastro de membros, um
aplicativo de projetos, um painel de pesquisa --- precisa \textbf{guardar dados
e recuperá-los depois}. Esta capacitação é sobre a ferramenta padrão da
indústria para isso: um banco de dados relacional (MySQL) operado pela
linguagem SQL.

Em duas horas não é possível formar um administrador de banco de dados. É
possível, sim, sair daqui \textbf{criando uma tabela, inserindo dados,
consultando com filtro e relacionando duas tabelas} --- exatamente o mínimo
necessário para o primeiro projeto real.

\section{Objetivos de aprendizagem}

Ao final da capacitação, o participante deve ser capaz de:

\begin{enumerate}
    \item Explicar o que é um SGBD e o que significa ``relacional''.
    \item Conectar-se a um servidor MySQL usando um cliente SQL.
    \item Criar um banco e uma tabela com tipos, chave primária e restrições
          adequadas.
    \item Executar as quatro operações do CRUD com \comando{INSERT},
          \comando{SELECT}, \comando{UPDATE} e \comando{DELETE}.
    \item Filtrar e ordenar resultados com \comando{WHERE}, \comando{ORDER BY}
          e \comando{LIMIT}.
    \item Relacionar duas tabelas com chave estrangeira e consultá-las com
          \comando{JOIN}.
\end{enumerate}

\section{O que não entra nesta capacitação}

Dito explicitamente para não gerar expectativa errada --- cada item destes é
uma capacitação inteira por si só:

\begin{itemize}
    \item Docker e containers (usados aqui apenas como meio de subir o ambiente).
    \item Normalização formal (1FN, 2FN, 3FN) e modelagem entidade-relacio\-namento completa.
    \item Índices, planos de execução e otimização de desempenho.
    \item Transações, \emph{views}, \emph{procedures}, \emph{triggers}, backup e replicação.
    \item Segurança, gestão de usuários e privilégios além do básico.
    \item Bancos não relacionais (MongoDB, Redis e afins).
\end{itemize}

\section{Roteiro de duas horas}
\label{sec:roteiro}

A Tabela~\ref{tab:roteiro} é o contrato de tempo da capacitação. Cada bloco tem
um tempo fechado. Se um bloco estourar, o corte sai do conteúdo marcado como
bônus (Seção~\ref{sec:agregacao}), nunca do bloco de \comando{JOIN}.

\begin{table}[H]
\centering
\caption{Distribuição dos 120 minutos da capacitação}
\label{tab:roteiro}
\footnotesize
\begin{tabular}{@{}clcc@{}}
\toprule
\textbf{\#} & \textbf{Bloco} & \textbf{Início} & \textbf{Duração} \\
\midrule
0 & Abertura, objetivos e escopo                       & 00:00 & 5 min  \\
1 & Conceitos-base: dado, banco, SGBD, relacional, SQL & 00:05 & 15 min \\
2 & Ambiente e cliente SQL: todos conectados           & 00:20 & 15 min \\
3 & DDL: banco, tipos, tabela e chave primária         & 00:35 & 20 min \\
-- & \emph{Pausa}                                      & 00:55 & 5 min  \\
4 & \comando{INSERT} e \comando{SELECT} com filtros    & 01:00 & 25 min \\
5 & \comando{UPDATE}, \comando{DELETE} e o \comando{WHERE} ausente & 01:25 & 10 min \\
6 & Chave estrangeira e \comando{JOIN}                 & 01:35 & 20 min \\
7 & Erros comuns, próximos passos e dúvidas            & 01:55 & 5 min  \\
\midrule
\multicolumn{3}{@{}l}{\textbf{Total}} & \textbf{120 min} \\
\bottomrule
\end{tabular}
\end{table}

\begin{dica}[Regra prática do instrutor]
Os blocos 3 a 6 são de prática. Fale no máximo um terço do tempo do bloco e
deixe dois terços para os participantes digitarem. Não avance enquanto a
maioria não tiver o resultado na tela.
\end{dica}

\section{Conceitos-base}
\label{sec:conceitos}

\subsection{Dado, informação e banco de dados}

\textbf{Dado} é um valor bruto: \comando{"Ana"}, \comando{2024-03-11},
\comando{7}. Sozinho, ele não diz nada. \textbf{Informação} é o dado dentro de
um contexto: \emph{``Ana ingressou na liga em 11/03/2024 e está no 7º
período''}. Um \textbf{banco de dados} é uma coleção organizada de dados,
guardada de forma persistente e estruturada para ser consultada.

\subsection{SGBD: o software servidor}

O \textbf{SGBD} (Sistema de Gerenciamento de Banco de Dados; em inglês,
\emph{DBMS}) é o software servidor que cuida desse banco. É ele quem de fato:

\begin{itemize}
    \item grava e lê os dados em disco;
    \item garante a integridade, não deixando entrar dado que viole as regras definidas;
    \item controla o acesso concorrente, permitindo que várias pessoas escrevam
          ao mesmo tempo sem corromper nada;
    \item controla as permissões dos usuários.
\end{itemize}

\begin{aviso}[Correção de um erro comum]
SQL \textbf{não} é a única forma de conversar com um SGBD. A comunicação real
acontece por um \textbf{protocolo de rede binário} --- o MySQL escuta na porta
3306. O SQL é a linguagem que trafega dentro desse protocolo: é a interface
padrão e dominante, mas não a única. O MySQL também expõe o \emph{X DevAPI /
Document Store}, uma API de CRUD sem SQL. E, quando se usa um ORM (Prisma,
Eloquent, Hibernate, SQLAlchemy), escrevem-se objetos --- é o ORM que traduz
aquilo para SQL antes de enviar.
\end{aviso}

\subsection{O modelo relacional}

O \textbf{modelo relacional}, proposto por Edgar F. Codd em 1970
\cite{codd1970}, é a ideia central: os dados são organizados em
\textbf{tabelas} (relações). Cada \textbf{linha} (registro ou tupla) é uma
ocorrência do mundo real; cada \textbf{coluna} (campo ou atributo) é uma
característica dela, com um \textbf{tipo} fixo. Tabelas se ligam a outras
tabelas por meio de \textbf{chaves} --- daí o nome ``relacional''.

\newpage

\begin{lstlisting}[style=base,caption={Anatomia de uma tabela},label={lst:anatomia}]
TABELA membro
+----+--------------+--------------------+---------+
| id | nome         | email              | periodo |  <- colunas (com tipo)
+----+--------------+--------------------+---------+
|  1 | Ana Souza    | ana@liga.edu.br    |       7 |  <- linha (registro)
|  2 | Bruno Lima   | bruno@liga.edu.br  |       3 |
+----+--------------+--------------------+---------+
  ^
  chave primaria: identifica cada linha de forma unica
\end{lstlisting}

\begin{conceito}[SGBD relacional]
Um \textbf{SGBD relacional} é um SGBD que implementa esse modelo --- é o que a
sigla \textbf{RDBMS} significa. MySQL, PostgreSQL, SQL Server, Oracle Database
e SQLite são todos RDBMS.
\end{conceito}
          % Introdução e conceitos-base
\chapter{MySQL}

\section{O que é o MySQL}

O MySQL é um \textbf{SGBD relacional} de código aberto, criado em 1995 e hoje
mantido pela Oracle. É o motor que armazena e gerencia os dados; a conversa com
ele acontece em SQL.

É provavelmente o banco relacional mais usado em aplicações web: WordPress,
boa parte dos sistemas em PHP e uma infinidade de APIs rodam sobre ele. A
versão usada nesta capacitação é a \textbf{8.4 LTS}, e o motor de armazenamento
padrão é o \textbf{InnoDB}.

\begin{conceito}[Por que o InnoDB importa]
É o InnoDB que dá suporte a \textbf{chaves estrangeiras} e a \textbf{transações}.
O motor antigo, MyISAM, não suporta chave estrangeira: ele aceita a declaração e
a ignora silenciosamente. Como todo o Capítulo~\ref{cap:relacionamentos} depende
de chave estrangeira, as tabelas deste guia são criadas com
\comando{ENGINE=InnoDB} explícito.
\end{conceito}

\section{MySQL e MariaDB: a confusão mais comum}

\begin{table}[H]
\centering
\caption{Comparação entre MySQL e MariaDB}
\label{tab:mysql-mariadb}
\footnotesize
\begin{tabularx}{\textwidth}{@{}lXX@{}}
\toprule
 & \textbf{MySQL} & \textbf{MariaDB} \\
\midrule
Origem & Original, de 1995 & \emph{Fork} de 2009, feito pelos criadores do
MySQL após a compra pela Oracle \\
\addlinespace
Licença da edição livre & \emph{Community Edition}: GPLv2, gratuita &
\emph{Server}: GPLv2, gratuita \\
\addlinespace
Edições pagas & Sim (\emph{Standard} e \emph{Enterprise}, da Oracle) &
Sim (suporte e produtos da MariaDB plc) \\
\addlinespace
Compatibilidade & \multicolumn{2}{X@{}}{Altíssima nos comandos básicos: tudo
desta capacitação funciona igual nos dois} \\
\bottomrule
\end{tabularx}
\end{table}

\newpage

\begin{aviso}[Correção de um erro comum]
Dizer que ``MariaDB é livre e MySQL é comercial'' é impreciso. \textbf{Ambos têm
uma edição livre em GPLv2 e ambos têm ofertas comerciais.} A diferença que
interessa agora é outra: várias distribuições Linux (Debian, Ubuntu, Fedora,
RHEL) empacotam o \textbf{MariaDB} como padrão nos repositórios oficiais, e o
XAMPP também entrega MariaDB. Para instalar o MySQL da Oracle
especificamente, é preciso adicionar o repositório oficial do MySQL ou usar
Docker.

Para o conteúdo desta capacitação, \textbf{tanto faz}: se o seu ambiente for
MariaDB, siga normalmente.
\end{aviso}

\section{Como executar o MySQL}

Existem três caminhos. Escolha apenas um.

\subsection{Opção A: pacote tudo-em-um}

Instala servidor web, PHP, banco e phpMyAdmin de uma vez só. É o caminho mais
fácil para quem está começando.

\begin{itemize}
    \item \href{https://www.apachefriends.org/}{\textbf{XAMPP}} --- Windows,
          Linux e macOS. É o mais popular. A sigla é \textbf{X}
          (\emph{cross-platform}) + \textbf{A}pache + \textbf{M}ariaDB +
          \textbf{P}HP + \textbf{P}erl. Em versões antigas o ``M'' era MySQL;
          hoje é MariaDB.
    \item \href{https://www.wampserver.com/}{\textbf{WampServer}} --- apenas
          Windows. Aí sim: \textbf{W}indows + \textbf{A}pache + \textbf{M}ySQL
          + \textbf{P}HP.
    \item \href{https://laragon.org/}{\textbf{Laragon}} --- apenas Windows,
          alternativa mais leve e moderna.
\end{itemize}

\begin{dica}[Conflito de portas]
Se o MySQL não subir, quase sempre há outro serviço ocupando a porta 3306
(outra instalação de MySQL, ou o SQL Server). A solução é mudar a porta no
\comando{my.ini} ou \comando{my.cnf} (\comando{port=3307}) e informar essa
porta no cliente.
\end{dica}

\subsection{Opção B: instalação nativa do servidor}

\begin{itemize}
    \item \textbf{Windows:} \emph{MySQL Installer for Windows}, no site oficial.
    \item \textbf{macOS:} \comando{brew install mysql} (Homebrew) ou o
          \comando{.dmg} oficial.
    \item \textbf{Linux (Debian/Ubuntu):} \comando{sudo apt install mysql-server}.
          Atenção: em várias distribuições esse pacote resolve para MariaDB.
          Para o MySQL da Oracle, adicione antes o repositório APT oficial.
\end{itemize}

Depois de instalar, execute \comando{sudo mysql\_secure\_installation} para
definir a senha do usuário \comando{root}.

\subsection{Opção C: Docker}

\begin{aviso}[Aviso de escopo]
Docker \textbf{não} faz parte desta capacitação. Está aqui apenas para
documentar como o ambiente da aula foi montado e para estimular quem quiser
pesquisar depois. Os arquivos prontos estão na pasta \comando{infrastructure/}
do repositório.
\end{aviso}

O arquivo de orquestração é um \comando{compose.yml}, que é \textbf{diferente
de um \comando{Dockerfile}}: o \comando{Dockerfile} descreve \emph{como
construir uma imagem}; o \comando{compose.yml} descreve \emph{quais serviços
prontos subir e como conectá-los}. Como aqui só usamos imagens prontas, só
existe \comando{compose.yml}.

\begin{bashcode}[caption={Subindo o ambiente},label={lst:compose-up}]
cd infrastructure
cp .env.example .env      # preencha MYSQL_ROOT_PASSWORD
docker compose up -d      # sobe MySQL + phpMyAdmin em segundo plano
docker compose ps         # confere se estao de pe e saudaveis
\end{bashcode}

\begin{bashcode}[caption={Derrubando o ambiente},label={lst:compose-down}]
docker compose down       # para os containers, PRESERVA os dados
docker compose down -v    # para e APAGA o volume (perde tudo)
\end{bashcode}

\begin{dica}[\comando{docker compose}, sem hífen]
O comando correto é \comando{docker compose} (versão 2, um plugin do Docker).
O antigo \comando{docker-compose} com hífen é a versão 1, descontinuada
desde 2023.
\end{dica}

A área dos alunos é preparada por um script que roda automaticamente na
primeira subida do container:

\begin{sqlcode}[caption={infrastructure/init/01\_create\_databases.sql (trecho)},label={lst:init}]
-- Banco de demonstracao do instrutor: so o root tem acesso.
CREATE DATABASE IF NOT EXISTS cap CHARACTER SET utf8mb4;

-- Area dos alunos. O GRANT abaixo e sobre um PADRAO de nome, nao sobre
-- um banco que ja exista: `_` e `%` sao curingas em GRANT, por isso o
-- underscore literal precisa do escape `\_`.
-- Resultado: 'aluno' pode criar e administrar liga_ana, liga_bruno, ...
-- e nenhum outro banco.
GRANT ALL PRIVILEGES ON `liga\_%`.* TO 'aluno'@'%';
FLUSH PRIVILEGES;
\end{sqlcode}

\begin{aviso}[Scripts de \emph{init} rodam uma única vez]
Os arquivos em \comando{init/} são executados apenas na primeira subida, com o
volume de dados vazio. Se você editar um deles depois, é preciso zerar o volume
para reaplicar: \comando{docker compose down -v \&\& docker compose up -d}.
\end{aviso}
          % MySQL
\chapter{Clientes SQL}

Um \textbf{cliente SQL} é o programa pelo qual \emph{você} conversa com o
servidor: ele abre a conexão, envia os comandos SQL e exibe o resultado em
tabela. O servidor (MySQL) e o cliente são programas separados --- podem,
inclusive, estar em máquinas diferentes.

\section{Panorama de clientes}

\begin{table}[H]
\centering
\caption{Clientes SQL mais comuns para MySQL}
\label{tab:clientes}
\footnotesize
\begin{tabularx}{\textwidth}{@{}lllX@{}}
\toprule
\textbf{Cliente} & \textbf{Tipo} & \textbf{Plataformas} & \textbf{Custo} \\
\midrule
\comando{mysql} (CLI) & Terminal & Win / Linux / macOS & Gratuito, vem com o servidor \\
MySQL Workbench       & Desktop  & Win / Linux / macOS & Gratuito (GPL) \\
DBeaver               & Desktop  & Win / Linux / macOS & \emph{Community} gratuito; PRO pago \\
HeidiSQL              & Desktop  & Windows             & Gratuito \\
phpMyAdmin            & \textbf{Web} & navegador       & Gratuito \\
TablePlus             & Desktop  & Win / Linux / macOS & Versão gratuita limitada; pago \\
DBVisualizer          & Desktop  & Win / Linux / macOS & Versão gratuita limitada; Pro pago \\
Navicat               & Desktop  & Win / Linux / macOS & \textbf{Pago} (apenas teste grátis) \\
\bottomrule
\end{tabularx}
\end{table}

\begin{aviso}[Correção de um erro comum]
Navicat, DBVisualizer e SQLPro Studio \textbf{não são gratuitos}: são produtos
comerciais, com versão de teste ou edição gratuita limitada. Não os recomende
como ``gratuitos''.
\end{aviso}

\section{phpMyAdmin}

O phpMyAdmin \textbf{não é um programa que se instala no sistema operacional}
como os outros da lista: é uma \textbf{aplicação web escrita em PHP} que roda
num servidor web e é acessada pelo \textbf{navegador}. É por isso que ele
aparece junto do Apache e do PHP em pacotes como o XAMPP.

Ele foi escolhido para esta capacitação por três motivos: não exige instalação
na máquina do participante (basta uma URL), mostra a estrutura do banco
visualmente enquanto se aprende, e tem uma aba \textbf{SQL} onde os comandos
são digitados de verdade.

\section{Conectando}
\label{sec:conectando}

Nesta capacitação \textbf{ninguém instala nada}. O instrutor roda o MySQL e o
phpMyAdmin na máquina dele e publica apenas o phpMyAdmin na internet por um
túnel. O acesso é por uma URL, no navegador.

\begin{table}[H]
\centering
\caption{Dados de acesso da turma}
\label{tab:acesso}
\footnotesize
\begin{tabularx}{\textwidth}{@{}lX@{}}
\toprule
\textbf{Campo} & \textbf{Valor} \\
\midrule
URL do phpMyAdmin & Ditada pelo instrutor no início da aula \\
Usuário           & \comando{aluno} \\
Senha             & \comando{aluno} \\
\bottomrule
\end{tabularx}
\end{table}

\begin{dica}[Tela de aviso do túnel]
Na primeira visita, o serviço de túnel mostra uma página de aviso antes do
site. Clique em \emph{``Visit Site''} uma vez e siga; ela não aparece de novo.
\end{dica}

\subsection{Cada participante tem o seu próprio banco}
\label{sec:banco-proprio}

O servidor é um só, compartilhado pela turma inteira. Se todos trabalhassem no
mesmo banco, o \comando{CREATE TABLE} do segundo participante já falharia com
\comando{ERROR 1050: Table already exists}, e um \comando{DELETE} de qualquer
um apagaria o dado de todos.

Por isso \textbf{cada participante cria o próprio banco}, chamado
\comando{liga\_} seguido do seu primeiro nome:

\begin{center}
\comando{liga\_ana} \quad \comando{liga\_bruno} \quad \comando{liga\_carla}
\quad \comando{liga\_diego}
\end{center}

\begin{aviso}[Regra do nome]
Apenas letras minúsculas, sem acento e sem espaço: \comando{liga\_joao}, nunca
\comando{liga\_João} nem \comando{liga\_joao silva}.

Onde este guia escrever \comando{liga\_SEUNOME}, use o seu. O login
\comando{aluno} só tem permissão em bancos que comecem com \comando{liga\_};
qualquer outro nome retorna \comando{ERROR 1044: Access denied}. Isso é
proposital: protege o banco de demonstração do instrutor (\comando{cap}) e os
dados internos do MySQL.
\end{aviso}

\subsection{Teste de sanidade}

Abra a aba \textbf{SQL} e execute:

\begin{sqlcode}
SELECT VERSION(), NOW();
\end{sqlcode}

Se voltou a versão do servidor e a data e hora, está tudo certo.
          % Clientes SQL
\chapter{SQL: fundamentos e definição de estruturas}
\label{cap:sql}

\textbf{SQL} (\emph{Structured Query Language}, ``linguagem de consulta
estruturada'') é a linguagem padronizada pela ANSI e pela ISO para definir e
manipular dados em bancos relacionais. Ela é \textbf{declarativa}: descreve-se
\emph{o que} se quer, não \emph{como} buscar. O SGBD decide o caminho.

Cada SGBD tem seu ``sotaque'' (dialeto), mas o núcleo é o mesmo --- o que se
aprende aqui vale, com ajustes pequenos, para PostgreSQL, SQL Server e SQLite.

\section{SQL não é um bloco só}

Os comandos SQL são agrupados por finalidade. Conhecer esse mapa evita muita
confusão.

\begin{table}[H]
\centering
\caption{Sublinguagens do SQL}
\label{tab:sublinguagens}
\footnotesize
\begin{tabularx}{\textwidth}{@{}llXl@{}}
\toprule
\textbf{Sigla} & \textbf{Nome} & \textbf{Para quê} & \textbf{Comandos} \\
\midrule
\textbf{DDL} & \emph{Data Definition Language}    & Define a \textbf{estrutura} & \comando{CREATE}, \comando{ALTER}, \comando{DROP} \\
\textbf{DML} & \emph{Data Manipulation Language}  & Manipula os \textbf{dados}  & \comando{INSERT}, \comando{UPDATE}, \comando{DELETE} \\
\textbf{DQL} & \emph{Data Query Language}         & \textbf{Consulta} os dados  & \comando{SELECT} \\
\textbf{DCL} & \emph{Data Control Language}       & Controla \textbf{permissões}& \comando{GRANT}, \comando{REVOKE} \\
\textbf{TCL} & \emph{Transaction Control Language}& Controla \textbf{transações}& \comando{COMMIT}, \comando{ROLLBACK} \\
\bottomrule
\end{tabularx}
\end{table}

\begin{dica}
Muitos autores tratam o \comando{SELECT} como parte do DML. Os dois usos
existem na literatura; o importante é entender a diferença entre \textbf{mexer
na estrutura} (DDL) e \textbf{mexer nos dados} (DML e DQL). Nesta capacitação
usamos DDL, DML e DQL.
\end{dica}

\section{A ordem de trabalho}

Como o MySQL é relacional, sempre existe uma estrutura antes do dado. A ordem
nunca muda:

\begin{lstlisting}[style=base,caption={Do banco vazio ao dado consultado},label={lst:ordem}]
1. CREATE DATABASE  ->  criar o banco          | DDL: a estrutura,
2. USE              ->  selecionar o banco     | feita uma vez
3. CREATE TABLE     ->  criar a(s) tabela(s)   |
                    v
4. INSERT           ->  inserir dados          | DML/DQL: o dia a dia,
5. SELECT           ->  consultar dados        | repetido o tempo todo
6. UPDATE           ->  alterar dados          |
7. DELETE           ->  remover dados          |
\end{lstlisting}

\begin{aviso}[Ponto que costuma passar batido]
Criar tabela e inserir dados não são o mesmo tipo de operação. Criar tabela é
DDL, feito uma vez no projeto. \comando{INSERT}, \comando{UPDATE} e
\comando{DELETE} são DML, executados milhares de vezes pela aplicação.
\end{aviso}

O cenário usado no restante deste guia é um cadastro de membros da liga
acadêmica, com duas tabelas: \comando{curso} e \comando{membro}.

\section{Criando o banco}
\label{sec:create-database}

\begin{aviso}
Troque \comando{SEUNOME} pelo seu primeiro nome, aqui e em todo o resto do
guia. Se você é a Ana, seu banco é \comando{liga\_ana}. Veja a
Seção~\ref{sec:banco-proprio}.
\end{aviso}

\begin{sqlcode}
CREATE DATABASE IF NOT EXISTS liga_SEUNOME
  CHARACTER SET utf8mb4
  COLLATE utf8mb4_0900_ai_ci;

USE liga_SEUNOME;
\end{sqlcode}

\begin{itemize}
    \item \comando{IF NOT EXISTS} evita erro se o banco já existir.
    \item \comando{utf8mb4} é o conjunto de caracteres que suporta
          \textbf{acentos e emojis} corretamente. É o padrão do MySQL 8, mas
          declará-lo explicitamente é boa prática: em servidores antigos o
          padrão era \comando{latin1}, que quebra a acentuação.
    \item \comando{USE liga\_SEUNOME;} define o banco corrente. \textbf{Sem
          isso, os comandos seguintes não sabem onde agir} e o servidor
          responde \comando{ERROR 1046: No database selected}. No phpMyAdmin,
          clicar no banco na barra lateral tem o mesmo efeito.
\end{itemize}

Comandos úteis de inspeção:

\begin{sqlcode}
SHOW DATABASES;              -- lista os bancos visiveis para voce
SHOW TABLES;                 -- lista as tabelas do banco corrente
DESCRIBE membro;             -- mostra as colunas e tipos de uma tabela
SHOW CREATE TABLE membro;    -- mostra o CREATE TABLE completo
\end{sqlcode}

\section{Tipos de dados}
\label{sec:tipos}

Cada coluna tem um \textbf{tipo}, e é ele que garante que a coluna
\comando{periodo} nunca receba \comando{"banana"}. Escolher bem o tipo é boa
parte da qualidade de uma tabela.

\begin{table}[H]
\centering
\caption{Tipos de dados mais usados no MySQL}
\label{tab:tipos}
\footnotesize
\begin{tabularx}{\textwidth}{@{}lXX@{}}
\toprule
\textbf{Tipo} & \textbf{Guarda} & \textbf{Quando usar} \\
\midrule
\comando{INT}          & Inteiro (cerca de $\pm$2,1 bilhões) & IDs, contadores, quantidades \\
\comando{TINYINT}      & Inteiro pequeno & Período, idade \\
\comando{BIGINT}       & Inteiro muito grande & IDs de sistemas de altíssimo volume \\
\comando{DECIMAL(p,s)} & Número exato com casas decimais & \textbf{Dinheiro}: \comando{DECIMAL(10,2)} \\
\comando{FLOAT} / \comando{DOUBLE} & Número aproximado & Medidas científicas. \textbf{Nunca para dinheiro} \\
\comando{VARCHAR(n)}   & Texto variável, até $n$ & Nome, e-mail, título: o caso mais comum \\
\comando{CHAR(n)}      & Texto de tamanho \textbf{fixo} & Siglas: \comando{CHAR(2)} para UF \\
\comando{TEXT}         & Texto longo (até 64 KB) & Descrições, observações \\
\comando{DATE}         & Data (\comando{AAAA-MM-DD}) & Data de nascimento, de ingresso \\
\comando{DATETIME}     & Data e hora & Agendamentos, registro de eventos \\
\comando{TIMESTAMP}    & Data e hora, sensível a fuso & \comando{criado\_em}, \comando{atualizado\_em} \\
\comando{BOOLEAN}      & Verdadeiro ou falso & Sinalizadores. No MySQL é apelido de \comando{TINYINT(1)} \\
\comando{ENUM(...)}    & Um valor de uma lista fixa & Status: \comando{ENUM('ativo','inativo')} \\
\bottomrule
\end{tabularx}
\end{table}

\begin{aviso}[Regra de ouro]
Dinheiro em \comando{DECIMAL}, \textbf{nunca} em \comando{FLOAT}.
\comando{FLOAT} é aproximado, e $0{,}1 + 0{,}2$ pode não dar exatamente
$0{,}3$.
\end{aviso}

\begin{conceito}[\comando{NULL} não é zero nem string vazia]
\comando{NULL} significa ``valor desconhecido ou não informado''. Por isso
\comando{WHERE curso\_id = NULL} \textbf{nunca} funciona: o correto é
\comando{WHERE curso\_id IS NULL}.
\end{conceito}

\section{Criando a tabela}
\label{sec:create-table}

\begin{sqlcode}
CREATE TABLE curso (
  id   INT AUTO_INCREMENT PRIMARY KEY,
  nome VARCHAR(80) NOT NULL UNIQUE
) ENGINE=InnoDB;
\end{sqlcode}

\begin{sqlcode}
CREATE TABLE membro (
  id            INT           AUTO_INCREMENT PRIMARY KEY,
  nome          VARCHAR(120)  NOT NULL,
  email         VARCHAR(160)  NOT NULL UNIQUE,
  periodo       TINYINT       NOT NULL,
  data_ingresso DATE          NOT NULL,
  ativo         BOOLEAN       NOT NULL DEFAULT TRUE,
  curso_id      INT           NULL,
  criado_em     TIMESTAMP     NOT NULL DEFAULT CURRENT_TIMESTAMP
) ENGINE=InnoDB;
\end{sqlcode}

As \textbf{restrições} (\emph{constraints}) são onde mora a qualidade do dado.

\begin{table}[H]
\centering
\caption{Restrições de coluna}
\label{tab:constraints}
\footnotesize
\begin{tabularx}{\textwidth}{@{}lX@{}}
\toprule
\textbf{Restrição} & \textbf{O que faz} \\
\midrule
\comando{PRIMARY KEY}    & Identifica cada linha de forma \textbf{única}. Não aceita \comando{NULL}, não repete. Uma por tabela \\
\comando{AUTO\_INCREMENT}& O MySQL gera o próximo número sozinho; esse campo não é informado no \comando{INSERT} \\
\comando{NOT NULL}       & O campo é obrigatório \\
\comando{UNIQUE}         & O valor não pode se repetir na coluna \\
\comando{DEFAULT x}      & Valor assumido quando o \comando{INSERT} não informa a coluna \\
\comando{FOREIGN KEY}    & Liga esta tabela a outra (Capítulo~\ref{cap:relacionamentos}) \\
\bottomrule
\end{tabularx}
\end{table}

\begin{conceito}[\comando{PRIMARY KEY} e \comando{UNIQUE}]
As duas impedem repetição. A diferença: só existe \textbf{uma} chave primária
por tabela e ela \textbf{não aceita \comando{NULL}}; já \comando{UNIQUE} pode
haver várias na mesma tabela, e ela aceita \comando{NULL}.
\end{conceito}

Alterar a estrutura depois também é DDL:

\begin{sqlcode}
ALTER TABLE membro ADD COLUMN telefone VARCHAR(20) NULL;
ALTER TABLE membro MODIFY COLUMN nome VARCHAR(150) NOT NULL;
ALTER TABLE membro DROP COLUMN telefone;

DROP TABLE membro;    -- apaga a tabela E todos os dados, sem desfazer
\end{sqlcode}
          % SQL: fundamentos e DDL
\chapter{CRUD: manipulando os dados}
\label{cap:crud}

\textbf{CRUD} é o acrônimo das quatro operações básicas sobre dados. Todo
sistema que você escrever na vida faz essas quatro coisas.

\begin{table}[H]
\centering
\caption{As quatro operações do CRUD}
\label{tab:crud}
\footnotesize
\begin{tabular}{@{}lll@{}}
\toprule
\textbf{CRUD} & \textbf{Comando SQL} & \textbf{Categoria} \\
\midrule
\textbf{C}reate & \comando{INSERT} & DML \\
\textbf{R}ead   & \comando{SELECT} & DQL \\
\textbf{U}pdate & \comando{UPDATE} & DML \\
\textbf{D}elete & \comando{DELETE} & DML \\
\bottomrule
\end{tabular}
\end{table}

\section{Create: \comando{INSERT}}

\begin{sqlcode}
-- Uma linha
INSERT INTO curso (nome) VALUES ('Engenharia de Software');

-- Varias linhas de uma vez (mais eficiente)
INSERT INTO curso (nome) VALUES
  ('Ciencia da Computacao'),
  ('Sistemas de Informacao'),
  ('Engenharia de Producao');

INSERT INTO membro (nome, email, periodo, data_ingresso, curso_id) VALUES
  ('Ana Souza',   'ana@liga.edu.br',   7, '2024-03-11', 1),
  ('Bruno Lima',  'bruno@liga.edu.br', 3, '2025-02-20', 2),
  ('Carla Dias',  'carla@liga.edu.br', 5, '2024-08-05', 1),
  ('Diego Rocha', 'diego@liga.edu.br', 9, '2023-09-14', 3),
  ('Elisa Prado', 'elisa@liga.edu.br', 1, '2026-03-02', NULL);
\end{sqlcode}

Repare que \comando{id}, \comando{ativo} e \comando{criado\_em} não foram
informados: são preenchidos pelo \comando{AUTO\_INCREMENT} e pelos
\comando{DEFAULT}.

\begin{itemize}
    \item Datas vão sempre no formato \comando{'AAAA-MM-DD'} e entre aspas.
    \item Texto entre aspas simples (\comando{'Ana'}); números sem aspas (\comando{7}).
    \item A ordem dos valores em \comando{VALUES} deve corresponder exatamente
          à ordem das colunas listadas.
\end{itemize}

\section{Read: \comando{SELECT}}

O comando mais usado da linguagem. A estrutura completa é:

\begin{sqlcode}
SELECT   colunas
FROM     tabela
WHERE    condicao          -- filtra LINHAS
ORDER BY coluna [ASC|DESC] -- ordena
LIMIT    n;                -- limita a quantidade
\end{sqlcode}

\begin{sqlcode}
SELECT * FROM membro;                    -- tudo (bom para explorar, ruim em producao)
SELECT nome, email FROM membro;          -- so as colunas que interessam
SELECT nome AS aluno, periodo AS sem FROM membro;  -- AS renomeia no resultado
\end{sqlcode}

\subsection{Filtrando com \comando{WHERE}}

\begin{sqlcode}
SELECT * FROM membro WHERE periodo > 5;
SELECT * FROM membro WHERE ativo = TRUE AND periodo >= 5;
SELECT * FROM membro WHERE periodo IN (1, 3, 5);
SELECT * FROM membro WHERE periodo BETWEEN 3 AND 7;    -- inclusivo nas 2 pontas
SELECT * FROM membro WHERE nome LIKE 'A%';             -- comeca com A
SELECT * FROM membro WHERE email LIKE '%@liga.edu.br'; -- termina com
SELECT * FROM membro WHERE nome LIKE '%ana%';          -- contem
SELECT * FROM membro WHERE curso_id IS NULL;           -- sem curso informado
SELECT * FROM membro WHERE data_ingresso >= '2025-01-01';
\end{sqlcode}

\begin{table}[H]
\centering
\caption{Operadores de filtro}
\label{tab:operadores}
\footnotesize
\begin{tabularx}{\textwidth}{@{}lX@{}}
\toprule
\textbf{Operador} & \textbf{Significado} \\
\midrule
\comando{=} \quad \comando{<>} (ou \comando{!=}) & Igual / diferente \\
\comando{>} \quad \comando{<} \quad \comando{>=} \quad \comando{<=} & Comparação \\
\comando{AND} \quad \comando{OR} \quad \comando{NOT} & Combinação lógica \\
\comando{IN (a, b, c)}      & Está na lista \\
\comando{BETWEEN a AND b}   & Está no intervalo, incluindo as pontas \\
\comando{LIKE}              & Padrão de texto: \comando{\%} é qualquer coisa, \comando{\_} é um caractere \\
\comando{IS NULL} / \comando{IS NOT NULL} & Testa ausência de valor \\
\bottomrule
\end{tabularx}
\end{table}

\subsection{Ordenando e limitando}

\begin{sqlcode}
SELECT nome, periodo FROM membro ORDER BY periodo DESC;
SELECT nome, periodo FROM membro ORDER BY periodo DESC, nome ASC;
SELECT nome FROM membro ORDER BY data_ingresso ASC LIMIT 3;  -- os 3 mais antigos
\end{sqlcode}

\comando{ASC} é crescente (padrão, pode ser omitido) e \comando{DESC} é
decrescente.

\section{Update: alterando}

\begin{sqlcode}
UPDATE membro
SET    periodo = 8
WHERE  id = 1;

UPDATE membro
SET    ativo = FALSE, periodo = 10
WHERE  email = 'diego@liga.edu.br';
\end{sqlcode}

\section{Delete: removendo}

Para praticar sem destruir os dados da aula, crie um registro descartável e
apague apenas ele:

\begin{sqlcode}
INSERT INTO membro (nome, email, periodo, data_ingresso)
VALUES ('Registro Teste', 'teste@liga.edu.br', 1, '2026-01-01');

SELECT * FROM membro WHERE email = 'teste@liga.edu.br';   -- confira ANTES
DELETE FROM membro WHERE email = 'teste@liga.edu.br';
\end{sqlcode}

\section{O erro mais caro de quem está começando}
\label{sec:where-ausente}

\begin{aviso}[Sem \comando{WHERE}, o comando vale para a tabela inteira]
\begin{sqlcode}
UPDATE membro SET periodo = 1;   -- zera o periodo de TODOS os membros
DELETE FROM membro;              -- apaga TODOS os membros da tabela
\end{sqlcode}
E não existe ``desfazer'' em SQL.
\end{aviso}

O hábito obrigatório é testar com \comando{SELECT} antes:

\begin{sqlcode}
-- 1) veja EXATAMENTE quais linhas serao afetadas
SELECT * FROM membro WHERE id = 5;

-- 2) so entao troque o SELECT * pelo DELETE ou UPDATE,
--    mantendo exatamente o mesmo WHERE
DELETE FROM membro WHERE id = 5;
\end{sqlcode}

\newpage

\begin{dica}[Rede de proteção, e a falta dela]
O MySQL Workbench tem o \emph{safe update mode}, que \textbf{bloqueia}
\comando{UPDATE} e \comando{DELETE} sem \comando{WHERE} em coluna-chave. O
phpMyAdmin \textbf{não} protege: nele, a disciplina é sua.

\comando{DELETE FROM tabela;} apaga as linhas uma a uma (é DML).
\comando{TRUNCATE TABLE tabela;} esvazia a tabela de forma muito mais rápida e
reinicia o \comando{AUTO\_INCREMENT}, mas é DDL e não pode ser desfeito por
\comando{ROLLBACK}.
\end{dica}
          % CRUD
\chapter{Relacionamentos: a parte relacional}
\label{cap:relacionamentos}

Até aqui usamos uma tabela só --- isso é uma planilha, não um banco relacional.
O que muda tudo é \textbf{relacionar} tabelas.

\section{O problema que a chave estrangeira resolve}

Se guardássemos o nome do curso dentro de \comando{membro}, teríamos
``Engenharia de Software'' repetido em dezenas de linhas. Aí alguém digita
``Eng. de Software'', outro digita ``engenharia de software'', e a consulta por
curso quebra. Corrigir o nome exigiria alterar todas as linhas.

A solução: o curso vira uma tabela própria, e \comando{membro} guarda apenas o
\comando{id} do curso.

\begin{lstlisting}[style=base,caption={Duas tabelas ligadas por chave},label={lst:fk}]
   curso                          membro
+----+---------------------+   +----+------------+----------+
| id | nome                |   | id | nome       | curso_id |
+----+---------------------+   +----+------------+----------+
|  1 | Eng. de Software    |<--|  1 | Ana Souza  |        1 |
|  2 | Ciencia da Comp.    |<--|  2 | Bruno Lima |        2 |
+----+---------------------+   |  3 | Carla Dias |        1 |
       ^                       +----+------------+----------+
  chave primaria (PK)                          chave estrangeira (FK)
\end{lstlisting}

\begin{conceito}[PK e FK]
A \textbf{chave primária (PK)} identifica a linha \emph{dentro} da própria
tabela. A \textbf{chave estrangeira (FK)} é uma coluna que \emph{aponta} para a
PK de outra tabela. Com ela declarada, o SGBD passa a \textbf{garantir} que
aquele valor existe do outro lado --- isso se chama \textbf{integridade
referencial}.
\end{conceito}

\section{Declarando a chave estrangeira}

\begin{sqlcode}
ALTER TABLE membro
  ADD CONSTRAINT fk_membro_curso
  FOREIGN KEY (curso_id) REFERENCES curso(id)
  ON DELETE SET NULL
  ON UPDATE CASCADE;
\end{sqlcode}

Com a chave estrangeira ativa, o comando abaixo passa a \textbf{falhar} --- e é
exatamente o que queremos:

\begin{sqlcode}
INSERT INTO membro (nome, email, periodo, data_ingresso, curso_id)
VALUES ('Fake', 'fake@liga.edu.br', 1, '2026-01-01', 999);
-- ERROR 1452: Cannot add or update a child row:
--             a foreign key constraint fails
\end{sqlcode}

\begin{table}[H]
\centering
\caption{Ações de \comando{ON DELETE} e \comando{ON UPDATE}}
\label{tab:fk-acoes}
\footnotesize
\begin{tabularx}{\textwidth}{@{}lX@{}}
\toprule
\textbf{Ação} & \textbf{Efeito ao apagar ou alterar o curso} \\
\midrule
\comando{RESTRICT} (padrão) & Impede a operação enquanto houver membro apontando para ele \\
\comando{SET NULL}          & Os membros ficam com \comando{curso\_id = NULL}; exige coluna que aceite \comando{NULL} \\
\comando{CASCADE}           & Propaga: apagar o curso apagaria os membros. \textbf{Use com muito cuidado} \\
\bottomrule
\end{tabularx}
\end{table}

\section{Cardinalidades}

\begin{itemize}
    \item \textbf{1:N (um para muitos)} --- o caso mais comum. Um curso tem
          vários membros; cada membro tem um curso. \textbf{A chave estrangeira
          fica no lado ``muitos''}, ou seja, em \comando{membro}. É o nosso caso.
    \item \textbf{1:1} --- chave estrangeira com \comando{UNIQUE}. Raro.
    \item \textbf{N:N (muitos para muitos)} --- por exemplo, um membro
          participa de vários projetos e um projeto tem vários membros.
          Resolve-se com uma \textbf{tabela intermediária}
          (\comando{membro\_projeto}) que guarda duas chaves estrangeiras.
          Conceito citado, não praticado nesta capacitação.
\end{itemize}

\section{\comando{JOIN}: consultando as duas tabelas juntas}

\comando{JOIN} combina linhas de tabelas diferentes usando a condição de
ligação declarada no \comando{ON}.

\begin{sqlcode}
-- INNER JOIN: so os membros QUE TEM curso
SELECT m.nome AS membro,
       c.nome AS curso,
       m.periodo
FROM   membro m
INNER  JOIN curso c ON m.curso_id = c.id
ORDER  BY c.nome, m.nome;
\end{sqlcode}

\newpage

\begin{sqlcode}
-- LEFT JOIN: TODOS os membros, tenham curso ou nao.
-- Elisa aparece com curso = NULL.
SELECT m.nome AS membro,
       c.nome AS curso
FROM   membro m
LEFT   JOIN curso c ON m.curso_id = c.id;
\end{sqlcode}

\begin{table}[H]
\centering
\caption{Tipos de \comando{JOIN}}
\label{tab:joins}
\footnotesize
\begin{tabularx}{\textwidth}{@{}lX@{}}
\toprule
\textbf{Tipo} & \textbf{Retorna} \\
\midrule
\comando{INNER JOIN} & Só as linhas com correspondência \textbf{nos dois lados} \\
\comando{LEFT JOIN}  & \textbf{Todas} as da tabela da esquerda, mais as correspondentes da direita (\comando{NULL} se não houver) \\
\comando{RIGHT JOIN} & O espelho do \comando{LEFT}; raro, basta inverter a ordem das tabelas \\
\bottomrule
\end{tabularx}
\end{table}

\begin{dica}[Apelidos e a armadilha do \comando{ON}]
\comando{membro m} e \comando{curso c} são \textbf{apelidos}
(\emph{aliases}). Com duas tabelas em jogo, \comando{m.nome} e \comando{c.nome}
deixam claro de qual tabela é cada coluna; sem isso o MySQL responde
\comando{ERROR 1052: Column 'nome' in field list is ambiguous}.

A condição do \comando{ON} é sempre \textbf{FK = PK}. Esquecer o \comando{ON}
produz um produto cartesiano: cada membro cruzado com cada curso.
\end{dica}

\section{Bônus: contar e agrupar}
\label{sec:agregacao}

\begin{dica}
Este conteúdo entra somente se o cronograma permitir
(Seção~\ref{sec:roteiro}). Caso contrário, fica como material de estudo.
\end{dica}

Funções de agregação resumem várias linhas em um valor: \comando{COUNT},
\comando{SUM}, \comando{AVG}, \comando{MIN} e \comando{MAX}.

\begin{sqlcode}
SELECT COUNT(*) FROM membro;      -- quantos membros existem
SELECT AVG(periodo) FROM membro;  -- periodo medio
SELECT MIN(data_ingresso), MAX(data_ingresso) FROM membro;
\end{sqlcode}

\comando{GROUP BY} aplica a agregação \textbf{por grupo}:

\begin{sqlcode}
SELECT c.nome        AS curso,
       COUNT(m.id)   AS total_membros
FROM   curso c
LEFT   JOIN membro m ON m.curso_id = c.id
GROUP  BY c.id, c.nome
ORDER  BY total_membros DESC;
\end{sqlcode}

\begin{conceito}[\comando{WHERE} e \comando{HAVING}]
\comando{WHERE} filtra \textbf{linhas antes} de agrupar; \comando{HAVING}
filtra \textbf{grupos depois} de agregar. Por exemplo:
\comando{... GROUP BY c.id, c.nome HAVING COUNT(m.id) >= 2;}
\end{conceito}
          % Relacionamentos e JOIN
\chapter{Erros comuns}
\label{cap:erros}

A tabela a seguir reúne as mensagens que mais aparecem para quem está
começando, com a causa e a correção de cada uma. Vale a pena consultá-la antes
de pedir ajuda: quase todo erro de aula está aqui.

\begin{longtable}{@{}p{4.6cm}p{4.4cm}p{4.6cm}@{}}
\caption{Erros frequentes, causas e correções}
\label{tab:erros} \\
\toprule
\textbf{Erro ou mensagem} & \textbf{Causa} & \textbf{Correção} \\
\midrule
\endfirsthead
\multicolumn{3}{@{}l}{\footnotesize\emph{Tabela~\ref{tab:erros} --- continuação}} \\
\toprule
\textbf{Erro ou mensagem} & \textbf{Causa} & \textbf{Correção} \\
\midrule
\endhead
\bottomrule
\endfoot
\footnotesize\comando{ERROR 1046: No database selected}
  & \footnotesize Faltou escolher o banco
  & \footnotesize \comando{USE liga\_SEUNOME;} ou clicar no banco na lateral \\
\addlinespace
\footnotesize\comando{ERROR 1044: Access denied to database}
  & \footnotesize Nome fora do padrão \comando{liga\_*}, ou tentativa de abrir o banco \comando{cap}
  & \footnotesize Use \comando{liga\_SEUNOME}: underscore, minúsculas, sem acento \\
\addlinespace
\footnotesize\comando{ERROR 1050: Table already exists}
  & \footnotesize O \comando{CREATE TABLE} foi executado duas vezes
  & \footnotesize \comando{DROP TABLE membro;} antes, ou rode o script do Apêndice~\ref{ap:script} inteiro \\
\addlinespace
\footnotesize\comando{ERROR 1062: Duplicate entry}
  & \footnotesize Violou \comando{PRIMARY KEY} ou \comando{UNIQUE}
  & \footnotesize O valor já existe; use outro ou faça \comando{UPDATE} \\
\addlinespace
\footnotesize\comando{ERROR 1452: foreign key constraint fails}
  & \footnotesize Chave estrangeira apontando para um \comando{id} inexistente
  & \footnotesize Insira primeiro na tabela-pai (\comando{curso}) \\
\addlinespace
\footnotesize\comando{ERROR 1452} ao apagar da tabela-pai
  & \footnotesize Há filhos referenciando aquela linha
  & \footnotesize Trate os filhos antes, ou defina \comando{ON DELETE} \\
\addlinespace
\footnotesize\comando{ERROR 1052: Column is ambiguous}
  & \footnotesize Duas tabelas com coluna de mesmo nome
  & \footnotesize Qualifique: \comando{m.nome}, \comando{c.nome} \\
\addlinespace
\footnotesize \comando{UPDATE} ou \comando{DELETE} afetou a tabela toda
  & \footnotesize \comando{WHERE} esquecido
  & \footnotesize Sempre testar com \comando{SELECT} antes (Seção~\ref{sec:where-ausente}) \\
\addlinespace
\footnotesize \comando{WHERE campo = NULL} não retorna nada
  & \footnotesize \comando{NULL} não é comparável com \comando{=}
  & \footnotesize \comando{WHERE campo IS NULL} \\
\addlinespace
\footnotesize Acentuação vira caractere estranho
  & \footnotesize Conjunto de caracteres errado (\comando{latin1})
  & \footnotesize Crie o banco com \comando{utf8mb4} \\
\addlinespace
\footnotesize Datas não filtram direito
  & \footnotesize Formato ou tipo errado
  & \footnotesize Use \comando{DATE}/\comando{DATETIME} e \comando{'AAAA-MM-DD'} \\
\addlinespace
\footnotesize Valores monetários erram centavos
  & \footnotesize Coluna declarada como \comando{FLOAT}
  & \footnotesize Use \comando{DECIMAL(10,2)} \\
\addlinespace
\footnotesize O comando não executa
  & \footnotesize Falta o ponto e vírgula final
  & \footnotesize Termine todo comando com \comando{;} \\
\end{longtable}
          % Erros comuns
\chapter{Conclusão e próximos passos}

Em duas horas percorremos o caminho completo de um banco relacional:
\textbf{entender o modelo}, \textbf{subir o servidor}, \textbf{conectar por um
cliente}, \textbf{definir a estrutura} (DDL), \textbf{manipular os dados}
(CRUD) e \textbf{relacionar tabelas} (chave estrangeira e \comando{JOIN}). Isso
é suficiente para modelar e operar o banco de um primeiro projeto da liga.

\section{O que estudar depois}

Nesta ordem:

\begin{enumerate}
    \item \textbf{Normalização (1FN, 2FN, 3FN)} --- como decidir quantas
          tabelas o seu projeto precisa \cite{elmasri2016}.
    \item \textbf{Modelagem entidade-relacionamento} --- desenhar o modelo
          antes de escrever SQL.
    \item \textbf{Relacionamento N:N} com tabela intermediária, e
          \comando{JOIN} de três ou mais tabelas.
    \item \textbf{Agregação avançada}: \comando{GROUP BY}, \comando{HAVING} e
          subconsultas.
    \item \textbf{Índices} --- por que uma consulta fica lenta e como
          acelerá-la (\comando{EXPLAIN}).
    \item \textbf{Transações} --- \comando{START TRANSACTION},
          \comando{COMMIT}, \comando{ROLLBACK} e o conceito ACID.
    \item \textbf{Integração com aplicação} --- conectar via linguagem
          (Python, PHP, Node) e conhecer os ORMs.
    \item \textbf{Backup e restauração} --- \comando{mysqldump}.
\end{enumerate}

\section{Onde praticar sem instalar nada}

\begin{itemize}
    \item \href{https://sqliteonline.com/}{SQLite Online} --- vários bancos direto no navegador
    \item \href{https://www.db-fiddle.com/}{DB Fiddle} --- MySQL e PostgreSQL no navegador
    \item \href{https://onecompiler.com/mysql/}{OneCompiler MySQL}
    \item \href{https://sqlzoo.net/}{SQLZoo} e \href{https://sqlbolt.com/}{SQLBolt} --- exercícios interativos guiados
    \item \href{https://www.hackerrank.com/domains/sql}{HackerRank SQL} --- desafios com correção automática
\end{itemize}

A documentação oficial do MySQL 8.4 \cite{mysql84} é a referência definitiva
para qualquer dúvida de sintaxe.

\section{Materiais desta capacitação}

\begin{itemize}
    \item \comando{docs/} --- este guia e a apresentação
    \item \comando{materials/liga.sql} --- o script completo da aula, pronto
          para colar no phpMyAdmin
    \item \comando{infrastructure/} --- \comando{compose.yml},
          \comando{.env.example} e os scripts de inicialização do banco
\end{itemize}
     % Conclusão e próximos passos

% --- Referências -------------------------------------------------------
\bibliographystyle{abntex2-alf}
\bibliography{pages/references}

% --- Elementos pós-textuais -------------------------------------------
\begin{apendicesenv}
\partapendices

% ----------------------------------------------------------------------
\chapter{Colinha de sintaxe}
\label{ap:colinha}

\begin{sqlcode}[caption={Estrutura (DDL)}]
CREATE DATABASE nome CHARACTER SET utf8mb4;
USE nome;
CREATE TABLE t (id INT AUTO_INCREMENT PRIMARY KEY, campo VARCHAR(100) NOT NULL);
ALTER TABLE t ADD COLUMN novo VARCHAR(50);
ALTER TABLE t MODIFY COLUMN campo VARCHAR(200) NOT NULL;
ALTER TABLE t DROP COLUMN novo;
DROP TABLE t;
\end{sqlcode}

\begin{sqlcode}[caption={Inspeção}]
SHOW DATABASES;   SHOW TABLES;   DESCRIBE t;   SHOW CREATE TABLE t;
\end{sqlcode}

\begin{sqlcode}[caption={Dados (DML e DQL)}]
INSERT INTO t (c1, c2) VALUES (v1, v2), (v3, v4);
SELECT c1, c2 FROM t WHERE cond ORDER BY c1 DESC LIMIT 10;
UPDATE t SET c1 = v WHERE id = 1;      -- SEMPRE com WHERE
DELETE FROM t WHERE id = 1;            -- SEMPRE com WHERE
\end{sqlcode}

\begin{sqlcode}[caption={Filtros}]
WHERE a = 1 AND b <> 2
WHERE a IN (1,2,3)          WHERE a BETWEEN 1 AND 10
WHERE nome LIKE 'A%'        WHERE campo IS NULL
\end{sqlcode}

\begin{sqlcode}[caption={Relacionamento}]
FOREIGN KEY (curso_id) REFERENCES curso(id) ON DELETE SET NULL;
SELECT m.nome, c.nome FROM membro m INNER JOIN curso c ON m.curso_id = c.id;
SELECT m.nome, c.nome FROM membro m LEFT  JOIN curso c ON m.curso_id = c.id;
\end{sqlcode}

\begin{sqlcode}[caption={Agregação}]
SELECT c.nome, COUNT(*) FROM curso c JOIN membro m ON m.curso_id = c.id
GROUP BY c.id, c.nome HAVING COUNT(*) >= 2;
\end{sqlcode}

% ----------------------------------------------------------------------
\chapter{Script completo da aula}
\label{ap:script}

Execute do início ao fim para reconstruir todo o cenário da capacitação. O
arquivo também está disponível no repositório, em
\comando{materials/liga.sql}.

\begin{aviso}[Antes de colar]
Substitua as três ocorrências de \comando{liga\_SEUNOME} pelo seu banco
--- \comando{liga\_ana}, \comando{liga\_joao} e assim por diante. Se você colar
sem trocar, vai criar um banco literalmente chamado \comando{liga\_seunome} e
disputá-lo com quem fizer o mesmo.
\end{aviso}

\lstinputlisting[style=sql]{../materials/liga.sql}

% ----------------------------------------------------------------------
\chapter{Exercícios}
\label{ap:exercicios}

Resolva no seu banco \comando{liga\_SEUNOME}, já construído pelo
Apêndice~\ref{ap:script}.

\begin{enumerate}
    \item Liste o nome e o e-mail de todos os membros \textbf{ativos}.
    \item Mostre os membros do 1º ao 5º período, do mais novo para o mais
          antigo na liga.
    \item Quantos membros ingressaram a partir de 2025?
    \item Cadastre um novo curso, \comando{'Engenharia Elétrica'}, e um membro
          nele.
    \item Corrija o e-mail da Carla para \comando{carla.dias@liga.edu.br}.
    \item Liste membro e nome do curso, incluindo quem está sem curso.
    \item Desative (\comando{ativo = FALSE}) todo membro que ingressou antes de
          \comando{2024-01-01}.
    \item Tente inserir um membro com \comando{curso\_id = 999}. Explique a
          mensagem de erro.
\end{enumerate}

\section{Respostas}

\begin{sqlcode}[caption={Exercícios 1 a 3}]
-- 1
SELECT nome, email FROM membro WHERE ativo = TRUE;

-- 2
SELECT * FROM membro
WHERE periodo BETWEEN 1 AND 5
ORDER BY data_ingresso DESC;

-- 3
SELECT COUNT(*) FROM membro WHERE data_ingresso >= '2025-01-01';
\end{sqlcode}

\begin{sqlcode}[caption={Exercícios 4 e 5}]
-- 4  (o curso PRIMEIRO: a chave estrangeira exige que o pai exista)
INSERT INTO curso (nome) VALUES ('Engenharia Eletrica');
INSERT INTO membro (nome, email, periodo, data_ingresso, curso_id)
VALUES ('Felipe Nunes', 'felipe@liga.edu.br', 2, '2026-03-10',
        (SELECT id FROM curso WHERE nome = 'Engenharia Eletrica'));

-- 5
SELECT * FROM membro WHERE nome = 'Carla Dias';           -- confira antes
UPDATE membro SET email = 'carla.dias@liga.edu.br' WHERE nome = 'Carla Dias';
\end{sqlcode}

\begin{sqlcode}[caption={Exercícios 6 a 8}]
-- 6  (LEFT JOIN, porque INNER JOIN esconderia quem nao tem curso)
SELECT m.nome AS membro, c.nome AS curso
FROM   membro m LEFT JOIN curso c ON m.curso_id = c.id;

-- 7
SELECT * FROM membro WHERE data_ingresso < '2024-01-01';  -- confira antes
UPDATE membro SET ativo = FALSE WHERE data_ingresso < '2024-01-01';

-- 8
INSERT INTO membro (nome, email, periodo, data_ingresso, curso_id)
VALUES ('Teste', 'teste@liga.edu.br', 1, '2026-01-01', 999);
-- ERROR 1452: a chave estrangeira exige que exista um curso com id = 999.
-- Como nao existe, o SGBD REJEITA a insercao. E a integridade referencial
-- protegendo o banco contra dado orfao.
\end{sqlcode}

\end{apendicesenv}


\end{document}
